\chapter{Study One: Experiencing Smart Learning Spaces}\label{ch:study-one}

\begin{synopsis}
\section{scoping human experience design in buildings}

The first part of the dissertation examines the forms of human experience that have already been addressed in research on indoor environments. Although many studies describe the improvement of occupant experience as an objective, they differ considerably in what they understand experience to include and how they seek to influence it.

The scoping review brings together research from HCI, architecture, building research, environmental psychology, and related fields. It examines the experiential aims pursued by this work, the design mechanisms used to address them, and the spatial or technological interventions developed. Its purpose is not only to classify previous research, but to establish the experiential territory within which HBI operates.

The review shows that human-experience design in buildings extends beyond dominant concerns with comfort, usability, and environmental performance. It includes emotional, bodily, aesthetic, social, democratic, inclusive, educational, and place-related forms of experience. It also identifies a broad repertoire of design approaches, including spatial, social, participatory, personalising, informational, playful, technological, and biophilic mechanisms.

At the same time, the review identifies emotional experience as a promising but under-examined way of approaching indoor complexity. Emotion is present across parts of the literature, yet it is often treated as an outcome to be measured rather than as a means of understanding how people relate to their surroundings. The review therefore points to the need for closer examination of the methodological potential of an emotional and affective lens.

This finding establishes the basis for the subsequent occupant studies. Rather than beginning with the assumption that comfort or satisfaction provides the most appropriate account of indoor experience, these studies examine what becomes visible when emotion is used as an entry point into occupants' situated accounts.
\end{synopsis}

\section{setting}
Describe the building, rooms, participants, and temporal context. Use margin figures for plans, routes, or small study artefacts.

\begin{marginfigure}
\fbox{\parbox{0.78\marginparwidth}{\centering\vspace{.7in}Map\vspace{.7in}}}
\captionof{figure}{A small margin figure is ideal for a floor-plan thumbnail or study route.}
\end{marginfigure}

\section{method}
Document recruitment, procedure, data collection, analysis, and ethical approval.

\section{findings}
\subsection{finding one}
State the finding as an analytic claim, then support it with data.

\subsection{finding two}
Use sidenotes for brief analytic qualifications rather than burying them in footnotes.

\section{discussion}
Connect the study back to the dissertation thesis.
